\documentclass[10pt,a4paper]{amsart}
\usepackage[margin=19mm]{geometry}
\usepackage{amsmath,amssymb,mathtools}
\usepackage[scr=rsfs,cal=euler]{mathalfa}
\usepackage{xfrac}
\usepackage[unicode,hidelinks,bookmarks=false]{hyperref}
\newcommand{\ie}{i.e.\ }
\newcommand{\cf}{cf.\ }
\newcommand{\resp}{respectively\ }
\newcommand{\Subsection}{\subsection}
\providecommand{\nobreakdash}{\nobreak\hbox{-}}
\setlength{\emergencystretch}{2em}
\multlinegap0pt
\usepackage{amssymb}
\usepackage[all]{xy}
\newtheorem{proposition}{Proposition}
\title{On the zero-classes of monoid semi-congruences}
\author{M. Hoefnagel, N. Martins-Ferreira, M. Sobral}
\date{Source: arXiv:2602.14015v1 (CC BY 4.0)}
\begin{document}
\maketitle

\noindent\emph{Source and licence.} On the zero-classes of monoid semi-congruences. Version arXiv:2602.14015v1 (CC BY 4.0), distributed by arXiv under the Creative Commons Attribution 4.0 International licence, http://creativecommons.org/licenses/by/4.0/, as recorded in the arXiv licence metadata for this exact version and shown on the abstract page https://arxiv.org/abs/2602.14015v1. Full licence notice: \texttt{licenses/arxiv-2602.14015-CC-BY-4.0.txt}.\par\smallskip

\noindent\textit{Editorial scope.} Counted unit: the article's proposition proposition (source0.tex lines 158--167). The article's complete original proof of it is delivered with it (source0.tex lines 169--182, one \texttt{proof} environment of 14 lines). No mathematics is written, repaired or re-derived: the statement and the proof are the article's own lines, in the article's own notation, and no retained line is substituted or re-flowed.  No further source-local context is needed: every cross-reference of every retained line resolves inside the two delivered spans.  Nothing was omitted from any retained span, so the package carries no omission record.  No retained line contains a citation, so the wrapper carries no bibliography and no bibliography entry.  No retained line contains a figure environment, an image-inclusion command or drawing code, so no image is delivered and the package declares no asset dependency.  Editorial formatting only, in addition: an AMS \texttt{amsart} preamble that restates the article's own declarations and macros used by the retained lines, namely the environments \texttt{proposition} on separate counters, together with the standard packages those lines need.  The retained environments print in this document as Proposition 1 in order of appearance, whereas the article prints the same items with the numbering of its own full document.  The complete native source of the article as distributed in the arXiv e-print for this exact version is bundled unchanged as \texttt{sources/194680/source0.tex}.
\begin{proposition}
 Let $M$ be a submonoid of a monoid $A$. Then $M$ is the zero-class of the reflexive relation $R_M$ if and only if the following condition holds:
\begin{quote}
(R) For every $x,y\in A$ and every $u\in M$,
\[
xy = 1 \;\Rightarrow\; xuy \in M.
\]
\end{quote}
If $M$ is a clot, then $M = [1]_{R_M}$, and the converse holds provided that the relation $R_M$ is internal.
\end{proposition}

\begin{proof}
We first show that $[1]_{R_M} \subseteq M$. Indeed, if $1 R_M u$, then for every $x,y\in A$ we have
\[
x1y = 1 \;\Rightarrow\; xuy \in M.
\]
In particular, taking $x=y=1$, we conclude that $u\in M$.

Conversely, $M \subseteq [1]_{R_M}$ if and only if condition (R) is satisfied. Consequently, under condition (R), we have
$M = [1]_{R_M}$.

Now assume that $M$ is a clot, that is, the zero-class of some internal reflexive relation $S$ on $A$. If $1Su$, then by reflexivity and compatibility of $S$ we have $xySxuy$. Thus, whenever $xy=1$, we obtain $1Sxuy$, that is, $xuy\in M$. Hence condition (R) holds and therefore $M=[1]_{R_M}$.

If $R_M$ is compatible, then the converse implication holds.
\end{proof}

\end{document}
